\section{Background}
\label{sec:background}

Containerization has fundamentally transformed cloud-native software deployment by offering
lightweight, scalable, and portable execution environments. Unlike virtual machines (VMs), which
encapsulate a full operating system per instance, containers share the host's OS kernel while
providing isolated user spaces via namespaces and control groups (cgroups). This architecture
reduces overhead and accelerates deployment, making containers ideal for microservices, CI/CD
pipelines, and elastic workloads. However, this shared kernel model introduces unique security
risks---particularly at the syscall interface, which governs access to core kernel
functionalities.

\subsection{System Call Security in Shared-Kernel Environments}
\label{subsec:syscall_security}

Syscalls act as the primary interface between user-space applications and kernel resources. In
containerized systems, they are critical to process execution, file operations, memory
management, and network communication. Unfortunately, their broad capabilities also make them a
prime target for attackers. The kernel exposes hundreds of syscalls, many of which are
unnecessary for typical containerized workloads, thereby expanding the attack
surface~\cite{ref11_shu_security}.

Attackers often exploit misused or vulnerable syscalls to execute privilege escalation, escape
containers, or tamper with the host environment. Such threats are exacerbated in multi-tenant
cloud deployments, where multiple containers share the same host kernel and thus the same syscall
namespace. Without rigorous syscall control, a compromise in one container can lead to lateral
movement or host takeover.

\subsection{Types of Syscall-Based Threats}
\label{subsec:syscall_threats}

System calls expose core kernel functions, making them prime targets in containerized
environments. Attackers exploit them to escalate privileges, escape isolation, or abuse
over-permissive profiles.

\textbf{Privilege Escalation:} Attackers may exploit kernel vulnerabilities, misconfigured Linux
capabilities, or incomplete seccomp profiles to gain elevated privileges from within a container.
Examples include the use of \texttt{setuid()}, \texttt{clone()}, or \texttt{execve()} in
specific sequences to hijack thread execution or bypass namespace restrictions.

\textbf{Container Escape:} Exploits like \mbox{CVE-2022-0847}~\cite{cve_2022_0847} (Dirty Pipe)
or \mbox{CVE-2019-5736}~\cite{cve_2019_5736} (Docker runc) allow malicious processes to bypass
container boundaries and interact with the host. These often involve crafting malicious syscall
sequences that tamper with read-only memory or escape namespace isolation.

\textbf{Attack Surface Abuse:} Many containers have overly broad access to syscalls due to
generic or default seccomp profiles. This over-permissiveness can be abused by sophisticated
adversaries to chain together syscalls into stealthy, multi-stage attacks.

\subsection{Limitations of Traditional Mitigation Approaches}
\label{subsec:traditional_limitations}

Security mechanisms such as AppArmor, SELinux, and Seccomp provide syscall filtering and policy
enforcement. However, they exhibit several limitations in modern cloud-native environments:

\begin{itemize}
  \item \textbf{Static Policy Inflexibility:} Static profiles do not adapt to new behaviors or
        threat vectors, often resulting in false positives (over-blocking) or false negatives
        (missed attacks).
  \item \textbf{Lack of Sequence Awareness:} Rule-based systems operate on individual syscall
        invocations without considering temporal dependencies, making them blind to multi-step
        attacks.
  \item \textbf{Scalability Bottlenecks:} As container density increases, static filters become
        harder to maintain and enforce efficiently at scale.
  \item \textbf{Manual Configuration Overhead:} Policies must be manually written, tested, and
        maintained for each container workload---an impractical burden in dynamic environments.
\end{itemize}

These shortcomings highlight the need for dynamic, intelligent, and context-aware syscall
security solutions that scale with modern infrastructure.

\subsection{The Promise of eBPF for Syscall Monitoring}
\label{subsec:ebpf}

eBPF is a kernel-native technology that allows safe, programmable instrumentation of kernel
events, including syscalls~\cite{ref20_bpf_security}. It enables dynamic attachment of monitoring
and filtering logic directly at syscall entry points (e.g., \texttt{sys\_enter\_*}) with minimal
overhead. Key benefits of eBPF in the syscall security domain include:

\begin{itemize}
  \item \textbf{Real-Time Enforcement:} Syscalls can be blocked or audited instantly without
        context switches or restarts.
  \item \textbf{Context-Aware Filtering:} Policies can account for namespaces, cgroup IDs,
        arguments, or syscall patterns.
  \item \textbf{High Scalability:} Efficient in-kernel maps and tracepoints allow enforcement
        over thousands of containers concurrently.
\end{itemize}

eBPF's programmable nature makes it an ideal foundation for implementing fine-grained, real-time
syscall security tailored to container workloads.

\subsection{Sequence-Based Anomaly Detection for Syscalls}
\label{subsec:sequence_anomaly}

Machine learning (ML) models, especially those analyzing syscall sequences, provide a powerful
complement to eBPF monitoring. Rather than relying solely on predefined rules, sequence-based
models can learn patterns of normal behavior and flag anomalous syscall windows that deviate from
expected workflows.

Variational Autoencoders (VAE)~\cite{vae_kingma} and Isolation
Forests~\cite{iforest_liu} (iForest) are particularly effective for modeling syscall sequences:

\begin{itemize}
  \item VAEs reconstruct normal behavior from compressed latent space representations. High
        reconstruction error indicates anomalies.
  \item Isolation Forests identify statistical outliers based on how easily a sequence can be
        separated from normal patterns.
  \item Ensembles of both improve robustness by combining reconstruction-based and tree-based
        anomaly scoring.
\end{itemize}

To support training and validation, public datasets such as the DongTing
dataset~\cite{ref26_dongtingdataset} offer syscall sequences labeled as normal or malicious.
DeSFAM uses DongTing exclusively, leveraging its 18,966 labeled sequences to train and calibrate
its anomaly detection models.

\subsection{Toward Adaptive, Scalable Syscall Security}
\label{subsec:adaptive_security}

The dynamic nature of container orchestration demands syscall filtering mechanisms that adapt to
workload behaviors, incorporate threat intelligence, and scale with minimal human intervention.

DeSFAM addresses these needs by combining:
\begin{itemize}
  \item Hybrid static-dynamic syscall profiling.
  \item Ensemble anomaly detection combining deep learning (VAE) and tree-based models (iForest).
  \item Risk-based policy enforcement using eBPF LSM hooks.
  \item CVE-matching and MITRE ATT\&CK-based risk scoring.
  \item Auto-hardening mechanisms to refine seccomp profiles at runtime.
\end{itemize}

This integrated design enables real-time detection, contextual blocking, and long-term syscall
surface reduction, setting the stage for resilient and scalable container security in
shared-kernel environments.
